% TODO make hw stylesheet
\documentclass[10pt]{article}
\usepackage[letterpaper,top=1in,left=1in,right=1in]{geometry}
\usepackage[dvipsnames,svgnames]{xcolor}
\usepackage[shortlabels]{enumitem}
\usepackage[framemethod=TikZ]{mdframed}
\usepackage{amsmath,amssymb,amsthm}
\usepackage[colorlinks]{hyperref}
\usepackage{multicol}
\usepackage{tabularx}
\usepackage{bbm}

\hypersetup{
    colorlinks=true,
    linkcolor=blue,
    filecolor=magenta,
    urlcolor=magenta,
    pdftitle={MAT 549 HW}
}

\usepackage{mathtools}
\usepackage{thmtools}
\usepackage[textsize=footnotesize]{todonotes}
\usepackage{titling}
\usepackage{cancel}

\reversemarginpar
\mdfdefinestyle{mdbluebox}{roundcorner=10pt,innerbottommargin=9pt,
linecolor=blue,backgroundcolor=TealBlue!5,}
\declaretheoremstyle[headfont=\sffamily\bfseries\color{MidnightBlue},
mdframed={style=mdbluebox},]{thmbluebox}

\mdfdefinestyle{mdbloobox}{frametitlefont=\bfseries,innerbottommargin=8pt,
nobreak=true,backgroundcolor=TealBlue!5,linecolor=blue,}
\declaretheoremstyle[headfont=\bfseries\color{MidnightBlue},
mdframed={style=mdbloobox},headpunct={\\[3pt]},postheadspace=0pt,]{thmbloobox}

\mdfdefinestyle{mdredbox}{frametitlefont=\bfseries,innerbottommargin=8pt,
nobreak=true,backgroundcolor=Salmon!5,linecolor=RawSienna,}
\declaretheoremstyle[headfont=\bfseries\color{RawSienna},
mdframed={style=mdredbox},headpunct={\\[3pt]},postheadspace=0pt,]{thmredbox}

\mdfdefinestyle{mdgreenbox}{linecolor=ForestGreen,backgroundcolor=ForestGreen!5,
linewidth=2pt,rightline=false,leftline=true,topline=false,bottomline=false,}
\declaretheoremstyle[headfont=\bfseries\sffamily\color{ForestGreen!70!black},
mdframed={style=mdgreenbox},headpunct={ --- },]{thmgreenbox}

\mdfdefinestyle{mdorangebox}{linecolor=RedOrange,backgroundcolor=RedOrange!5,
linewidth=2pt,rightline=false,leftline=true,topline=false,bottomline=false,}
\declaretheoremstyle[headfont=\bfseries\sffamily\color{RedOrange!70!black},
mdframed={style=mdorangebox},headpunct={ --- },]{thmorangebox}

\mdfdefinestyle{mdblackbox}{linecolor=black,backgroundcolor=RedViolet!5!gray!5,
linewidth=3pt,rightline=false,leftline=true,topline=false,bottomline=false,}
\declaretheoremstyle[mdframed={style=mdblackbox}]{thmblackbox}

\declaretheoremstyle[spaceabove=6pt,spacebelow=6pt]{boldhead}
\declaretheoremstyle[spaceabove=6pt,spacebelow=6pt,headfont=\normalfont\itshape]{italichead}

\declaretheorem[style=boldhead,name=Problem]{problem}
\declaretheorem[style=boldhead,name=Problem,numbered=no]{problem*}
\declaretheorem[style=boldhead,name=Setup]{setup} % for config geo
\declaretheorem[style=boldhead,name=Example,sibling=problem]{example}
\declaretheorem[style=boldhead,name=$\bigstar$ Problem,sibling=problem]{niceprob}
\declaretheorem[style=boldhead,name=Theorem,sibling=problem]{theorem}
\declaretheorem[style=boldhead,name=Corollary,sibling=theorem]{corollary}
\declaretheorem[style=boldhead,name=Proposition,sibling=theorem]{prop}
\declaretheorem[style=boldhead,name=Lemma,numberwithin=problem]{lemma}
\declaretheorem[style=boldhead,name=Theorem,numbered=no]{theorem*}
\declaretheorem[style=boldhead,name=Lemma,numbered=no]{lemma*}
\declaretheorem[style=boldhead,name=Claim,numberwithin=problem]{claim}
\declaretheorem[style=boldhead,name=Claim,numbered=no]{claim*}
\declaretheorem[style=boldhead,name=Remark,numberwithin=problem]{remark}
\declaretheorem[style=boldhead,name=Remark,numbered=no]{remark*}
\declaretheorem[style=boldhead,name=Definition,sibling=theorem]{definition}
\declaretheorem[style=boldhead,name=Definition,numbered=no]{definition*}

\providecommand{\ol}{\overline}
\providecommand{\ceil}[1]{\left\lceil #1 \right\rceil}
\providecommand{\floor}[1]{\left\lfloor #1 \right\rfloor}
\newcommand{\dd}{\mathrm{d}}
\providecommand{\ul}{\underline}
\providecommand{\eps}{\varepsilon}
\providecommand{\half}{\frac{1}{2}}
\providecommand{\CC}{\mathbb C}
\providecommand{\FF}{\mathbb F}
\providecommand{\II}{\mathbb I}
\providecommand{\NN}{\mathbb N}
\providecommand{\QQ}{\mathbb Q}
\providecommand{\RR}{\mathbb R}
\providecommand{\ZZ}{\mathbb Z}
\providecommand{\ind}{\mathbbm 1}
\providecommand{\dg}{^\circ}
\providecommand{\ii}{\item}
\providecommand{\setdiff}{\smallsetminus}
\providecommand{\alert}[1]{{\sffamily\textbf{\textcolor{blue}{#1}}}}
\providecommand{\inv}{^{-1}}
\providecommand{\mb}[1]{\mathbf{#1}}
\newcommand{\what}{\widehat}

\newenvironment{soln}{\begin{proof}[Solution]}{\end{proof}}

\providecommand{\pow}{\operatorname{Pow}}
\providecommand{\vocab}[1]{{\textbf{\textcolor{ForestGreen}{#1}}}}
\providecommand{\lcm}{\operatorname{lcm}}
\providecommand{\abs}[1]{\left\lvert #1\right\rvert}
\providecommand{\smabs}[1]{\lvert #1\rvert}
\providecommand{\innerprod}[1]{\langle\!\langle #1\rangle\!\rangle}
\providecommand{\norm}[1]{\left\| #1\right\|}
\providecommand{\smnorm}[1]{\| #1\|}
\providecommand{\gr}{\operatorname{gr}}

\usepackage{mathrsfs}
\providecommand{\tanvec}[1]{\frac{\partial}{\partial #1}}

% place title at top right
\setlength{\droptitle}{-8ex}
\pretitle{\begin{flushright}\Large\bfseries}
\posttitle{\par\end{flushright}}
\preauthor{\begin{flushright}}
\postauthor{\end{flushright}}
\predate{\begin{flushright}}
\postdate{\end{flushright}}

\title{MAT 549 HW04}
\author{\large Aditya Pahuja}
\date{\large Due 03/10}
\begin{document}
\maketitle

\begin{problem}
    [12-8]
    Let $M$ be a smooth $n$-manifold,
    and let $A$ be a smooth covariant $k$-tensor field on $M$.
    If $(U,(x_i))$ and $(\tilde U,\tilde x_j)$ are
    overlapping smooth charts on $M$, we can write
    \begin{equation*}
	A
	= \sum_{1\le i_1,\dots,i_k\le n}A_{i_1,\dots,i_k}
	\dd x_{i_1}\otimes\cdots\otimes\dd x_{i_k}
	= \sum_{1\le j_1,\dots,j_k\le n}\tilde A_{j_1,\dots,j_k}
	\dd \tilde x_{j_1}\otimes\cdots\otimes\dd \tilde x_{j_k}.
    \end{equation*}
    Compute a transformation law analogous to (11.7)
    expressing the component functions $A_{i_1,\dots,i_k}$
    in terms of $\tilde A_{j_1,\dots,j_k}$.
\end{problem}

\begin{soln}
    By equation (11.7) in Lee,
    \begin{equation*}
	\dd x_i =
	\sum_{j=1}^n\frac{\partial \tilde x_j}{\partial x_i}\dd\tilde x_j.
    \end{equation*}
    By definition of tensors, the smooth function commutes
    with $\otimes$, so
    \begin{equation*}
	\sum_{1\le i_1,\dots,i_k\le n}A_{i_1,\dots,i_k}
	\dd x_{i_1}\otimes\cdots\otimes\dd x_{i_k}
	= 
	\sum_{1\le i_1,\dots,i_k\le n}\tilde A_{i_1,\dots,i_k}
	\left(\sum_{j_1=1}^n\frac{\partial \tilde x_{j_1}}
	{\partial x_{i_1}}\dd\tilde x_{j_1}\right)
	\otimes\cdots\otimes
	\left(\sum_{j_k=1}^n\frac{\partial \tilde x_{j_k}}
	{\partial x_{i_k}}\dd\tilde x_{j_k}\right).
    \end{equation*}
    This means that the coefficient of a particular simple tensor
    $\dd\tilde x_{\ell_1}\otimes\cdots\otimes\dd\tilde x_{\ell_k}$ is
    \begin{equation*}
	A_{\ell_1,\dots,\ell_k} =
	\boxed{\sum_{1\le i_1,\dots,i_k\le n} \tilde A_{i_1,\dots,i_k}
	\frac{\partial\tilde x_{\ell_1}}{\partial x_{i_1}}
	\cdots\frac{\partial\tilde x_{\ell_k}}{\partial x_{i_k}}}.\qedhere
    \end{equation*}
\end{soln}

\begin{problem}
    [14-1]
    Show that covectors $\omega_1$, \dots, $\omega_k$
    on a finite-dimensional vector space $V$ are linearly dependent
    if and only if $\omega_1\wedge\cdots\wedge\omega_k = 0$.
\end{problem}

\begin{soln}
    If $\omega_1$, \dots, $\omega_k$ are linearly dependent,
    there exist real numbers $c_1$, \dots, $c_k$, not all zero, such that
    $c_1\omega_1 + \cdots + c_k\omega_k = 0$.
    Since reordering the vectors does not change whether
    their wedge product is nonzero (only affecting the sign),
    we can assume that $c_1\ne 0$, so that
    $\omega_1 = \frac {c_2}{c_1}\omega_2 \wedge
    \cdots \wedge \frac{c_k}{c_1}\omega_k$. Then
    \begin{equation*}
        \omega_1\wedge\omega_2\wedge\cdots\wedge\omega_k
	= \sum_{j=2}^k \frac{c_j}{c_1}\omega_j\wedge\omega_2\cdots\omega_k
	= 0
    \end{equation*}
    where each summand is zero because $\omega_j$ appears twice in the wedge.

    Conversely suppose that $\omega_1$, \dots, $\omega_k$
    are linearly independent, so that we can find covectors
    $\omega_{k+1}$, \dots, $\omega_n$ to make the list
    $\omega_1$, \dots, $\omega_n$ a basis of $V$. Then
    $\omega_1\wedge\cdots\wedge\omega_n$ is nonzero
    (since otherwise $\Lambda^nV^*$ would be zero-dimensional,
    because any element of $\Lambda^nV^*$ can be written
    as a scalar multiple of $\omega_1\wedge\cdots\wedge\omega_n$).
    This implies $\omega_1\wedge\cdots\wedge\omega_k$ is nonzero too,
    because the wedging with zero always yields zero.
\end{soln}

\begin{problem}
    [14-5]
    Let $M$ be a smooth $n$-manifold with or without boundary
    and let $(\omega_1,\dots,\omega_k)$ be an ordered $k$-tuple
    of smooth $1$-forms on an open subset $U\subseteq M$
    such that $(\omega_1|_p,\dots,\omega_k|_{p})$
    is linearly independent for each $p\in U$.
    Given smooth $1$-forms $\alpha_1$, \dots, $\alpha_k$ on $U$ such that
    \begin{equation*}
	\sum_{i=1}^k \alpha_i\wedge\omega_i = 0,
    \end{equation*}
    show that each $\alpha_i$ can be written as
    a linear combination of the $\omega_i$ with smooth coefficients.
\end{problem}

\begin{soln}
    Let $p\in M$ be arbitrary,
    and let $U$ be a small enough open set around $p$
    such that $(\omega_1,\dots,\omega_k)$ can be extended to
    a list $(\omega_1,\dots,\omega_n)$ for which
    $(\omega_1|_p,\dots,\omega_n|_p)$ is a basis for $T^*_pM$
    for each $p\in U$. Then, for each $1\le i\le k$ and $1\le j\le n$
    there exist $A_{i,j}\in C^\infty(U)$ such that
    \begin{equation*}
	\alpha_i = \sum_{j=1}^nA_{i,j}\omega_j.
    \end{equation*}
    For a particular $1\le i\le k$,
    \begin{equation*}
	0 = \left(\sum_{i'=1}^n\alpha_{i'}\wedge\omega_{i'}\right)
	\wedge(\omega_1\wedge\cdots\wedge\omega_{i-1}
	\wedge\omega_{i+1}\wedge\cdots\wedge\omega_k)
	= (-1)^{i-1}\alpha_i\wedge(\omega_1\wedge\cdots\wedge\omega_k)
    \end{equation*}
    This implies that $(\alpha_i,\omega_1,\dots,\omega_k)$
    are linearly dependent at each $p\in M$,
    so, in $U$, the coefficients $A_{\ell,j}$ of $\omega_\ell$
    for $k+1\le\ell\le n$ in the coordinate representation
    of $\alpha_i$ are all zero (noting that the $\omega_j$
    are nonzero because they span $T^*_pM$ at each $p\in U$).
    Thus $\alpha_i = \sum_{j=1}^k A_{i,j}\omega_j$ in $U$.

    If $\alpha_i = \sum_{j=1}^k A_{i,j}\omega_j$ in $U$
    and $\alpha_i = \sum_{j=1}^k B_{i,j}\omega_j$ in $V$
    for some open subsets $U,V\subseteq M$,
    then by the linear independence of the $\omega_j|_p$
    we can conclude that $A_{i,j} = B_{i,j}$ on $U\cap V$.
    Thus the coefficients $A_{i,j}^{U_p}$ in a neighborhood $U_p$ of $p\in M$
    agree with each other inside intersections of sets
    in the open cover $\{U_p\}$, so they define smooth functions
    $A_{i,j}\colon M\to\RR$ globally.
\end{soln}

\begin{problem}
    [14-6]
    Define a $2$-form $\omega$ on $\RR^3$ by
    \begin{equation*}
        \omega = x\dd y\wedge\dd z + y\dd z \wedge\dd x + z\dd x\wedge\dd y.
    \end{equation*}
    \begin{enumerate}[(a)]
        \ii Compute $\omega$ in spherical coordinates
	$(\rho,\psi,\theta)$ defined by $(x,y,z)
	= (\rho\sin\varphi\cos\theta,
	\rho\sin\varphi\sin\theta,
	\rho\cos\varphi)$.
	\ii Compute $\dd\omega$ in both Cartesian and spherical coordinates
	and verify that both expressions represent the same $3$-form.
	\ii Compute the pullback $\iota^*_{S^2}\omega$ to $S^2$,
	using coordinates $(\varphi,\theta)$ on
	the open subset where these coordinates are defined.
	\ii Show that $\iota^*_{S^2}\omega$ is nowhere zero.
    \end{enumerate}
\end{problem}

\begin{soln}
    (a) We have
    \begin{align*}
	\dd x
	&= \sin\varphi\cos\theta\dd\rho
	+ \rho\cos\varphi\cos\theta\dd\varphi
	- \rho\sin\varphi\sin\theta\dd\theta\\
	\dd y
	&= \sin\varphi\sin\theta\dd\rho
	+ \rho\cos\varphi\sin\theta\dd\varphi
	+ \rho\sin\varphi\cos\theta\dd\theta\\
	\dd z
	&= \cos\varphi\dd\rho - \rho\sin\varphi\dd\varphi\\
	\dd x\wedge\dd y
	&= \rho^2\cos\varphi\sin\varphi\dd\varphi\wedge\dd\theta
	+ \rho\sin^2\varphi\dd\rho\wedge\dd\theta\\
	\dd y\wedge\dd z
	&= -\rho\sin\theta\dd\rho\wedge\dd\varphi
	+ \rho^2\sin^2\varphi\cos\theta\dd\varphi\wedge\dd\theta
	- \rho\sin\varphi\cos\varphi\cos\theta\dd\rho\wedge\dd\theta
	\\
	\dd z\wedge\dd x
	&= \rho\cos\theta\dd\rho\wedge\dd\varphi
	+ \rho^2\sin^2\varphi\sin\theta\dd\varphi\wedge\dd\theta
	- \rho\sin\varphi\cos\varphi\sin\theta\dd\rho\wedge\dd\theta
    \end{align*}
    so
    \begin{equation*}
        \omega = \rho^3\sin\varphi\dd\varphi\wedge\dd\theta.
    \end{equation*}

    (b) In Cartesian coordinates,
    \begin{equation*}
        \dd\omega = \dd x\wedge\dd y\wedge\dd z
	+ \dd y\wedge \dd z \wedge\dd x
	+ \dd z\wedge\dd x\wedge\dd y
	= 3\dd x\wedge\dd y\wedge\dd z,
    \end{equation*}
    noting that the cyclic permutations of $x$, $y$, and $z$ are even.
    In spherical coordinates, we get
    \begin{equation*}
        \dd\omega = 3\rho^2\sin\varphi\dd\rho\wedge\dd\varphi\wedge\dd\theta
    \end{equation*}
    which is equal to the expression in Cartesian coordinates
    by Corollary 14.21 in Lee:
    \begin{equation*}
	\dd\rho\wedge\dd\varphi\wedge\dd\theta
        = \det \begin{pmatrix}
	    \sin\theta\cos\theta & \rho\cos\varphi\cos\theta &
	    -\rho\sin\varphi\sin\theta\\
	    \sin\varphi\sin\theta & \rho\cos\varphi\sin\theta &
	    \rho\sin\varphi\cos\theta\\
	    \cos\varphi & -\rho\sin\varphi & 0
        \end{pmatrix}\dd x\wedge\dd y\wedge\dd z
	= \rho^2\sin\varphi\dd\rho\wedge\dd\varphi\wedge\dd\theta.
    \end{equation*}

    (c) In spherical coordinates, $\iota_{S^2}(\varphi,\theta)
    = (1,\varphi,\theta)$, so the pullback of $\omega$ by the inclusion is
    \begin{equation*}
	\iota^*_{S^2}\omega = \sin\varphi\dd(\varphi\circ\iota_{S^2})\wedge
	\dd(\theta\circ\iota_{S^2}) = \sin\varphi\dd\varphi\wedge\dd\theta.
    \end{equation*}

    (d) The coordinate system $(\varphi,\theta)$
    is defined away from the poles, and in this region
    $0 < \varphi < \pi$, so
    $\omega = \sin\varphi\dd\varphi\wedge\dd\theta$ is nonzero
    away from the poles.
    A similar argument applied after swapping the $x$ and $z$ axes
    shows that $\omega$ is nonzero away from the
    intersections of $S^2$ with the $x$-axis, so $\iota^*_{S^2}\omega$
    is nowhere zero.
\end{soln}











\end{document}
